% EEE3096S Practical 3 report. Compile from this directory: latexmk -pdf report.tex
% The scope captures full12.jpg and coarse6.jpg sit beside this file.
\documentclass[11pt,a4paper]{article}
\usepackage{lmodern}
\usepackage[margin=20mm]{geometry}
\usepackage{amsmath}
\usepackage{booktabs}
\usepackage{tabularx}
\usepackage{graphicx}
\usepackage[hidelinks]{hyperref}
\setlength{\parindent}{0pt}
\setlength{\parskip}{5pt}
\renewcommand{\arraystretch}{1.4}
\newcommand{\blank}[1][27mm]{\rule{#1}{0.4pt}}
\newcommand{\capture}[1]{%
  \includegraphics[width=\linewidth,height=0.26\textheight,keepaspectratio]{#1}}

\begin{document}
\begin{center}
  {\Large\bfseries EEE3096S Practical 3 --- Voltage Mirror}
\end{center}
Names/student numbers: CHNJAC009 & SBYSIB014\par
Date: 23 October 2026\qquad Build configuration: -O0

\section*{1. Measurements and predictions}
Measured analogue reference $V_R$: 3.2946~V\par
DMM range/resolution: 10~V DC range, 5 significant figures (0.1~mV resolution)\par
Nominal ADC/DAC interval $V_R/4096$: $3.2946/4096 = 0.8043$~mV\par
Predicted coarse output step $V_R/64$: $3.2946/64 = 51.48$~mV

{\small
\begin{tabularx}{\linewidth}{@{}l l X r r X X@{}}
\toprule
Input target & Mode & Measured PA5 (V) & $C$ & $D$ & Predicted PA4 (V) & Measured PA4 (V) \\
\midrule
% Predicted PA4 = D x 3.2946 / 4096, using the D recorded in each row.
$\sim$0.800 V & FULL12  & 0.79955 & 992 & 992 & 0.7979 & 0.80163\\
$\sim$0.800 V & COARSE6 & 0.79955 & 992& 960 & 0.7722 & 0.77765\\
$\sim$1.65 V & FULL12  & 1.6528 & 2045 & 2045 & 1.6449 & 1.6536\\
$\sim$1.65 V & COARSE6 & 1.6528 & 2053 & 2048 & 1.6473 & 1.6536\\
$\sim$2.500 V & FULL12  & 2.5040 & 3113 & 3113& 2.5039 & 2.4753\\
$\sim$2.500 V & COARSE6 & 2.5040 & 3116 & 3072 & 2.4710 & 2.4753\\
\bottomrule
\end{tabularx}}

Generator: 1.996 Hz; 1.77 Vpp; 1.787 V offset\par
Verified PA5 minimum/maximum: 0.9 / 2.67 V\par
Measured adjacent coarse step: 51.25 mV (average over six steps)\par
Measured minus predicted step: $51.25 - 51.48 = -0.23$~mV\par
Percentage difference $100(\Delta V_{\rm measured}-\Delta V_{\rm predicted})/
\Delta V_{\rm predicted}$: $-0.44$~\%

\newpage
\section*{2. Scope captures}

\begin{center}
\capture{full12.jpg}
\end{center}
\textbf{Capture 1 --- FULL12.} CH1 = PA5 input, CH2 = PA4 output, both
500~mV/div, DC coupled; 50~ms/div (zoomed window of the 100~ms/div main
timebase). The PA4 trace lies on top of PA5 over the whole cycle
(PA5: 0.90--2.67~V, 1.998~Hz).

\begin{center}
\capture{coarse6.jpg}
\end{center}
\textbf{Capture 2 --- COARSE6.} CH1 = PA5 input (yellow), CH2 = PA4 output
(green), both 50~mV/div, DC coupled; 10~ms/div, rising ramp. PA4 is a
staircase of plateaus below the input ramp; the average spacing of adjacent
settled plateaus over six steps is 51.25~mV.

\newpage
\section*{3. Explanation}
\textbf{Mapping.} FULL12 copies the full 12-bit code, $D = C$, so all 4096 commands
are available. COARSE6 limits resolution by clearing the six least significant bits,
$D = C \mathbin{\&} \mathtt{0b111111000000} = 64\lfloor C/64\rfloor$ This gives us $4096/64 = 64$ distinct commands
(0, 64, \dots, 4032) spaced 64 codes apart, or \ $V_R/64 =
51.48$~mV at the output. The largest is 4032 since the next multiple of 64, 4096, needs a 13th bit. Because
bits are cleared, COARSE6 always rounds down ($D \le C$), so the output sits
0--63~LSB (0--51~mV) below the input, as the staircase in Capture~2 shows.
The ADC and DAC both remain 12-bit; only the software mapping removes the
six bits, giving an effective 6-bit resolution.

\textbf{Measured steps.} The measured step of 51.25~mV, which we averaged over six
adjacent plateaus, is 0.23~mV ($-0.44$\,\%) below the prediction
$51.48$~mV. This difference is within the expected measurement
uncertainty, which comes mainly from noise on the plateaus, cursor placement
and the scope's limited vertical resolution, so the measured step agrees with
$V_R/64$.

\textbf{Resolution and accuracy.} The FULL12 trace in Capture~1 only
\emph{looks} smooth. The DAC still outputs discrete codes, updated every
0.5~ms, so PA4 is a staircase of steps of about 3.5~mV on this triangle. These
steps are too small to see at 500~mV/div, so a smooth trace shows only that
the steps are below the scope's resolution, not that the output is continuous.

\emph{Resolution} is the smallest output step: 0.80~mV in FULL12 and
51.48~mV in COARSE6. \emph{Accuracy} is how close the output is to the true
value, and is set by the converters' offset, gain and linearity errors and by
noise, not by the number of bits. Our FULL12 outputs at 0.8~V and 1.65~V are
3.7--8.7~mV or 5--11~LSB from their predicted voltages, an error of about
0.5\,\%; that is to say, the conversion is fairly accurate.

\textbf{One DC result (0.8~V, FULL12).} With $C = 992$, the ADC estimate
$C V_R/4096 = 0.7979$~V is 1.6~mV ($\approx 2$~LSB) below the measured input
0.7996~V. With $D = C = 992$, the nominal DAC output $D V_R/4096 = 0.7979$~V
is 3.7~mV ($\approx 5$~LSB) below the measured output 0.8016~V, so the output
is 2.1~mV above the input. Copying the code is exact only digitally: the ADC
has its own offset, gain and linearity error, the DAC has independent errors, and noise makes $C$ vary
by a few LSB.
\end{document}
